\chapter{Trabalhos Relacionados}
\label{cap:trabalhos-relacionados}

Este capítulo posiciona a proposta em relação a três corpos de literatura que raramente se encontram em um mesmo trabalho: os modelos de difusão de informação e contágio, as redes neurais em grafos heterogêneos e temporais, e a previsão de popularidade e sucesso musical. A revisão é narrativa e orientada por relevância ao escopo da dissertação (difusão musical em charts brasileiros de streaming), e não pretende ser sistemática. O objetivo é evidenciar o que cada linha já oferece, onde elas se limitam e qual lacuna a combinação proposta pretende preencher. Enquanto o Capítulo~\ref{cap:fundamentacao} apresentou os conceitos, aqui o foco é o que trabalhos anteriores fizeram com eles.

%=====================================================

\section{Difusão de informação e modelos de contágio}

O estudo da difusão de inovações e comportamentos tem raízes anteriores às plataformas digitais. \citet{rogers2003diffusion} consolidou a teoria da difusão de inovações, descrevendo como uma novidade se espalha por uma população ao longo do tempo segundo curvas de adoção tipicamente sigmoides. Na mesma tradição, o modelo de \citet{bass1969} formalizou a adoção de produtos por meio de coeficientes de inovação e imitação, oferecendo uma das primeiras descrições paramétricas e ajustáveis de curvas de difusão, um antecedente direto da leitura epidemiológica adotada nesta dissertação.

Quando a difusão passa a ser mediada por uma rede, a estrutura das conexões torna-se determinante. \citet{granovetter1978} propôs modelos de limiar para comportamento coletivo, nos quais um indivíduo adota um comportamento apenas quando uma fração suficiente de seus contatos já o adotou. \citet{watts2002cascades} mostrou, com um modelo simples de cascatas em redes aleatórias, que grandes cascatas globais são possíveis mas raras, e dependem fortemente da conectividade e da distribuição de limiares. Essa linha culmina, do ponto de vista algorítmico, no trabalho de \citet{kempe2003ic}, que formalizou os modelos Independent Cascade e Linear Threshold e o problema de maximização de influência. Esses modelos tornam explícito que a topologia da rede condiciona o alcance da difusão, mas operam sobre grafos homogêneos e tipicamente sobre um único tipo de relação.

No contexto de conteúdo online, dois trabalhos ajudam a caracterizar a difusão para além do volume agregado. \citet{goel2016virality} introduziram a noção de \emph{virality estrutural}, distinguindo conteúdos que se espalham por difusão de muitos passos daqueles impulsionados por uma única transmissão massiva. \citet{zhao2015seismic}, por sua vez, modelaram a popularidade de publicações por meio de processos pontuais auto-excitados, prevendo o tamanho final de uma cascata a partir de sua dinâmica inicial. Esses trabalhos aproximam difusão e previsão de popularidade, mas permanecem centrados em uma única modalidade de interação (compartilhamento ou retweet) e não exploram a heterogeneidade de entidades.

A aplicação de modelos compartimentais a fenômenos culturais é mais recente e diretamente relevante a esta proposta. \citet{oliveira2025brasnam} ajustaram um modelo SIR a curvas de viralidade e sucesso de músicas em charts brasileiros, estimando parâmetros de adoção e remoção por faixa e estabelecendo o baseline epidemiológico que esta dissertação reproduz. Extensões posteriores incorporaram comportamento multi-onda a essas curvas \citep{oliveira2025asonam}. O ganho dessa abordagem é a interpretabilidade: poucos parâmetros resumem a trajetória de cada música. A limitação, contudo, é estrutural: o ajuste é feito música a música, de forma isolada, sem representar que artistas, gêneros e trajetórias semelhantes condicionam a difusão.

%=====================================================

\section{Redes neurais em grafos heterogêneos e temporais}

As redes neurais em grafos oferecem um caminho para incorporar exatamente a estrutura relacional que os modelos de contágio agregados deixam de fora. As arquiteturas fundamentais (GCN \citep{kipf2017gcn}, GraphSAGE \citep{hamilton2017graphsage} e GAT \citep{velickovic2018gat}) estabeleceram o paradigma de passagem de mensagens para grafos homogêneos, diferindo na forma de normalizar, amostrar ou ponderar vizinhos. Esses modelos, no entanto, tratam todos os nós e arestas como de um único tipo.

Para domínios com entidades e relações semanticamente distintas, surgiram as GNNs heterogêneas. A R-GCN \citep{schlichtkrull2018rgcn} introduziu parâmetros específicos por tipo de relação, permitindo modelar grafos de conhecimento com múltiplas arestas. A HAN \citep{wang2019han} combinou atenção em nível de nó e de \emph{meta-caminho}, e o \emph{Heterogeneous Graph Transformer} \citep{hu2020hgt} generalizou a atenção por tipo de nó e aresta. Essas arquiteturas justificam a escolha, nesta dissertação, de um codificador que trate \texttt{performs}, \texttt{has\_genre}, \texttt{cotrajectory} e \texttt{cooccurs} com parâmetros próprios, em vez de colapsá-las em um grafo homogêneo.

A dimensão temporal foi incorporada por duas famílias de abordagens, conforme discutido na fundamentação. A primeira discretiza o tempo em \emph{snapshots} e combina uma GNN com um módulo recorrente: a \emph{Graph Convolutional Recurrent Network} \citep{seo2018gcrn} encadeia convoluções em grafo e células recorrentes, e a EvolveGCN \citep{pareja2020evolvegcn} faz os próprios pesos da GNN evoluírem no tempo. A DySAT \citep{sankar2020dysat} aplica autoatenção tanto na estrutura quanto na sequência de snapshots. A segunda família preserva os \emph{timestamps} individuais e atualiza representações por evento: a TGAT \citep{xu2020tgat} introduz uma codificação temporal funcional sobre a vizinhança, e a TGN \citep{rossi2020tgn} mantém memória por nó atualizada a cada interação. Levantamentos da área sistematizam esse espaço de projeto \citep{kazemi2020dynamicsurvey}. A arquitetura proposta nesta dissertação situa-se claramente na primeira família (\emph{snapshots} semanais codificados por uma GNN heterogênea e agregados por uma rede recorrente), escolha justificada pelo grão semanal das relações do grafo e pelo custo computacional. As variantes contínuas, como a TGN, permanecem como alternativa para cenários de granularidade mais fina.

Apesar da maturidade dessas arquiteturas, sua aplicação conjunta (heterogeneidade \emph{e} temporalidade) ao problema específico de difusão musical é escassa, e praticamente inexistente sobre dados brasileiros de streaming.

%=====================================================

\section{Previsão de popularidade e sucesso musical}

A previsão de sucesso musical é estudada há quase duas décadas sob o rótulo de \emph{hit song science}. \citet{dhanaraj2005hit} foram dos primeiros a prever hits a partir de atributos de áudio e letra. Pouco depois, \citet{pachet2008hitscience} contestaram esses resultados, argumentando que atributos acústicos isolados têm poder preditivo limitado, um debate que permanece pertinente e que recomenda cautela ao tratar features de conteúdo como suficientes. Trabalhos mais recentes retomaram o problema com aprendizado profundo: \citet{yang2017revisiting} aplicaram redes convolucionais a espectrogramas para prever popularidade, e \citet{interiano2018musical} analisaram tendências e a previsibilidade do sucesso em décadas de músicas dos charts, encontrando previsibilidade modesta porém significativa a partir de características de áudio.

Do ponto de vista de dados, o MGD+ \citep{seufitelli2023dsw} consolidou metadados, atributos acústicos e séries de charts em escala global, viabilizando estudos reprodutíveis de difusão e popularidade. É sobre esse ecossistema, combinado a rankings diários de Top 200 e Viral 50, que esta dissertação se apoia.

O traço comum dessa literatura é o predomínio de representações tabulares ou de conteúdo: cada música é descrita por seu próprio vetor de atributos, e o sucesso é previsto de forma quase independente entre faixas. Raramente a estrutura relacional (coautoria, compartilhamento de gênero, similaridade de trajetória) ou a dinâmica temporal explícita da difusão são incorporadas ao modelo preditivo.

%=====================================================

\section{Posicionamento e lacuna}

As três linhas revisadas são complementares e, em conjunto, delimitam a contribuição desta dissertação. Os modelos de difusão e contágio oferecem interpretabilidade e capturam a dinâmica temporal, mas operam em nível populacional e ignoram a heterogeneidade de entidades; é o caso do baseline SIR aplicado a charts brasileiros \citep{oliveira2025brasnam}. As GNNs heterogêneas e temporais oferecem exatamente a representação relacional e dinâmica ausente nos modelos de contágio, mas foram pouco aplicadas à difusão musical e quase nunca a dados brasileiros. A literatura de previsão de popularidade musical, por fim, domina o uso de atributos de conteúdo, mas tende a tratar as músicas isoladamente, sem grafo nem difusão explícita.

A lacuna que emerge é a ausência de um trabalho que (i) represente a difusão musical como um grafo heterogêneo temporal de músicas, artistas e gêneros; (ii) compare diretamente esse modelo a um baseline epidemiológico interpretável e já consolidado na literatura brasileira; e (iii) o faça sob um protocolo de avaliação que respeite a ordem temporal de observação das relações. É precisamente nesse ponto que esta dissertação se insere: ao confrontar um modelo populacional, GNNs estáticas e uma GNN heterogênea temporal em um mesmo pipeline experimental sobre charts brasileiros, busca-se quantificar o ganho marginal da estrutura relacional e da temporalidade na previsibilidade da difusão musical.
